\documentclass[11pt,a4paper]{article}
\usepackage[T1]{fontenc}
\usepackage[utf8]{inputenc}
\usepackage{lmodern}
\usepackage{mathpazo}
\usepackage{amsmath,amssymb,amsthm,mathtools}
\usepackage[margin=27mm,headheight=14pt,headsep=9mm,footskip=12mm]{geometry}
\usepackage{microtype,fancyhdr,booktabs,longtable,array,enumitem,xcolor,seqsplit}
\usepackage[hyphens]{url}
\usepackage{hyperref}
\definecolor{paperteal}{HTML}{216C68}
\definecolor{paperrule}{gray}{0.72}
\hypersetup{colorlinks=true,linkcolor=paperteal,citecolor=paperteal,urlcolor=paperteal,
 pdftitle={A finite Fourier-budget criterion for global Navier--Stokes regularity on the three-torus},
 pdfsubject={Weak Fourierman B: ten-step derivation and formal verification},
 pdfkeywords={Navier-Stokes, Fourier budget, continuation, critical Sobolev data, Lean}}
\urlstyle{tt}
\Urlmuskip=0mu plus 2mu
\setlength{\parindent}{1.2em}
\setlength{\parskip}{3pt plus 1pt}
\setlength{\emergencystretch}{2em}
\linespread{1.04}
\pagestyle{fancy}
\fancyhf{}
\fancyhead[L]{\small Periodic Navier--Stokes equations}
\fancyhead[R]{\small Finite Fourier budgets}
\fancyfoot[L]{\small Preprint draft\enspace$\cdot$\enspace19 November 2026}
\fancyfoot[R]{\small\thepage}
\renewcommand{\headrulewidth}{0.35pt}
\renewcommand{\headrule}{\hbox to\headwidth{\color{paperrule}\leaders\hrule height \headrulewidth\hfill}}
\renewcommand{\footrulewidth}{0pt}
\numberwithin{equation}{section}
\theoremstyle{plain}
\newtheorem{theorem}{Theorem}[section]
\newtheorem{proposition}[theorem]{Proposition}
\newtheorem{lemma}[theorem]{Lemma}
\newtheorem{corollary}[theorem]{Corollary}
\theoremstyle{definition}
\newtheorem{definition}[theorem]{Definition}
\theoremstyle{remark}
\newtheorem{remark}[theorem]{Remark}
\newcommand{\T}{\mathbb T^3}
\newcommand{\R}{\mathbb R}
\newcommand{\Z}{\mathbb Z}
\newcommand{\D}{\mathcal D}
\newcommand{\HH}[1]{\dot H^{#1}}
\newcommand{\norm}[2]{\lVert #1\rVert_{#2}}
\newcommand{\pair}[2]{\langle #1,#2\rangle}
\newcommand{\Tmax}{T_{\max}}
\newcommand{\tq}{\tau_Q}
\newcommand{\Budget}{\mathcal B}
\newcommand{\Psol}{\mathbb P_\sigma}
\newcommand{\decl}[1]{\begingroup\small\nolinkurl{#1}\endgroup}
\newcommand{\hash}[1]{{\small\ttfamily\seqsplit{#1}}}
\newcommand{\step}[2]{\par\medskip\noindent\textbf{Step #1. #2}\par\smallskip}
\newcommand{\sourcecommit}{fdb515d8bde35ec8021bbc26f360aca8dfa2855c}

\begin{document}
\begin{center}
{\LARGE\bfseries A finite Fourier-budget criterion\par
for global Navier--Stokes regularity\par
on the three-torus\par}
\vspace{7pt}
{\large\color{paperteal} From critical Sobolev smallness to same-trajectory continuation\par}
\end{center}
\vspace{3pt}
\noindent\textbf{Abstract.}
For the unforced Navier--Stokes equations on the unit three-torus, we present
a finite-horizon Fourier criterion for global smooth continuation of the
unique maximal classical solution. The hypothesis is finiteness of the
squared full Fourier enstrophy integrated up to the smaller of the maximal
lifetime and the explicit time $\tq=\norm{a}{2}^4/(8\pi^2\nu^5)$.
An enstrophy inequality and Gronwall give a uniform $H^1$ bound; a common
local restart time then forces strict passage beyond $\tq$.
An inherited QFB theorem supplies a smooth tail of the same original
Leray--Hopf trajectory, and normalized-pressure agreement permits a smooth
global splice. We organize the argument into ten steps and connect it to the
small-$H^{1/2}$ theorem of \cite{small}, obtaining the explicit budget
$\norm{a}{\HH{1/2}}^4/[64\pi^4(\nu-P\norm{a}{\HH{1/2}})]$ under strict
critical smallness. The exact WFB endpoint has matching external Lean
compilation and theorem-level axiom-audit evidence. The finite-budget
hypothesis is retained throughout; it is not established for every smooth datum.

\medskip
\noindent\textbf{Keywords.} Periodic Navier--Stokes equations; Fourier enstrophy;
finite-horizon continuation; critical Sobolev smallness; Leray--Hopf solutions;
formal verification.

\medskip
\noindent\textbf{Verification status.}
The exact conditional WFB endpoint and its contextual obligations are
externally verified at the source revision recorded in Appendix~\ref{app:evidence}.
This paper is a mathematical exposition of that endpoint and a reconstruction
of its connection to \cite{small}. The arbitrary-admissible-constant
small-data argument is not claimed to be a newly formalized proof.

\section{Equation and scope}
We consider
\begin{equation}\label{eq:nse}
 \partial_tu+(u\cdot\nabla)u+\nabla p=\nu\Delta u,
 \qquad \nabla\cdot u=0,\qquad u(0)=a
\end{equation}
on $\T=\R^3/\Z^3$, with volume one, viscosity $\nu>0$ and zero forcing.
The initial datum belongs to
\[
 \D=C^\infty_{\sigma,0}(\T;\R^3)
 =\left\{a\in C^\infty(\T;\R^3):\nabla\cdot a=0,\ \int_\T a=0\right\}.
\]
Velocity and pressure are periodic, and pressure is fixed by
$\int_\T p(t,x)\,dx=0$. Smoothness at $t=0$ is understood on the closed
forward spacetime domain, with right time derivatives there.

The starting publication \cite{small} proves global regularity for small
critical Sobolev data. Its analytic argument contains two distinct components:
a smallness-dependent coercive estimate that generates a quartic gradient
budget, and a continuation argument that uses such a budget. The latter
component does not require the datum itself to be small.

The present criterion, called \emph{Weak\_Fourierman\_B} (WFB) in the formal
development, uses a budget only on the open interval ending at
$\min\{\Tmax,\tq\}$. The additional QFB result supplies eventual smoothness
after $\tq$ for the same original weak trajectory. The proof must therefore
establish a \emph{strict} overlap between the classical interval and the
regular tail. Merely reaching the cutoff is insufficient.

The analytical dependence is
\[
\begin{gathered}
\text{finite Fourier budget}\ \Longrightarrow\ \text{uniform }H^1\text{ bound}
\ \Longrightarrow\ \tq<\Tmax,\\
\text{strict overlap}+\text{same-trajectory QFB tail}
\ \Longrightarrow\ \text{global smooth continuation}.
\end{gathered}
\]
Critical smallness gives one explicit sufficient condition for the first
hypothesis. The paper does not prove that this hypothesis holds for all of
$\D$, nor does it claim a result for general rough data through the WFB
statement. The rough-$H^{1/2}$ trace and compactness theory in \cite{small}
remains a separate result.

\section{Normalization and the main theorem}\label{sec:main}
\subsection{Fourier coefficients and physical quantities}
For a periodic real vector field, set
\begin{equation}\label{eq:fourier}
 \widehat u(k)=\int_\T u(x)e^{-2\pi i k\cdot x}\,dx,
 \qquad k\in\Z^3,
 \qquad
 \norm{u}{\HH{s}}^2=\sum_{k\ne0}(2\pi|k|)^{2s}|\widehat u(k)|^2.
\end{equation}
Coefficient norms are Euclidean norms on $\mathbb C^3$.
Write $\Lambda=(-\Delta)^{1/2}$ and use the paper's notation
\begin{equation}\label{eq:energies}
\begin{aligned}
 E(u)&=\norm{u}{\HH{1/2}}^2,&
 Y(u)&=\norm{\nabla u}{L^2(\mathrm{HS})}^2,\\
 Z(u)&=\norm{u}{\HH{3/2}}^2,&
 W(u)&=\norm{\Delta u}{2}^2.
\end{aligned}
\end{equation}
Thus $E$ is a critical Sobolev quantity, not the kinetic energy
$\frac12\norm{u}{2}^2$. The full Fourier enstrophy is
\begin{equation}\label{eq:Q}
 Q_u(t)=\sum_{k\in\Z^3\setminus\{0\}}|k|^2|\widehat u(t,k)|^2,
 \qquad Y(u(t))=4\pi^2Q_u(t).
\end{equation}
All modes are present. No finite-mode truncation or numerical certificate
is substituted for this sum.

\subsection{Maximal trajectories and the budget}
A classical initial front of duration $T>0$ is a smooth normalized solution
of \eqref{eq:nse} on $[0,T)$. Local existence and uniqueness construct a
unique maximal pair $(u,p)$ with lifetime $\Tmax\in(0,\infty]$; every shorter
front is its restriction, and it agrees with every other classical solution
with the same initial datum on their common interval. Section~\ref{sec:local}
recalls the construction independently of any budget hypothesis.

Define
\begin{equation}\label{eq:cutoff}
 \tq(a,\nu)=\frac{\norm{a}{2}^4}{8\pi^2\nu^5},
 \qquad L=\min\{\Tmax,\tq(a,\nu)\},
 \qquad \Budget_L(u)=\int_{(0,L)}Q_u(t)^2\,dt.
\end{equation}
The integral is a nonnegative extended Lebesgue integral. Since $\tq$ is
finite, $L$ is finite even when $\Tmax=\infty$. If $L=0$, the open interval
is empty and the budget is zero. In the formal statement the minimum is
taken in the extended nonnegative reals \emph{before} conversion to a real
endpoint; Appendix~\ref{app:semantics} records the precise expression.

\begin{theorem}[Weak\_Fourierman\_B]\label{thm:wfb}
Let $\nu>0$ and $a\in\D$. There exists a unique maximal classical trajectory
$(u,p)$ with $\Tmax\in(0,\infty]$. If
\begin{equation}\label{eq:mainbudget}
 \Budget_{\min\{\Tmax,\tq(a,\nu)\}}(u)<\infty,
\end{equation}
then $\Tmax=\infty$. There is a global smooth solution pair $(U,\Pi)$ of
\eqref{eq:nse}, with zero velocity mean and normalized pressure, such that
\begin{equation}\label{eq:preservation}
 U(t)=u(t),\qquad \Pi(t)=p(t),\qquad 0\le t<\Tmax.
\end{equation}
In particular, the conclusion preserves the entire original trajectory.
\end{theorem}

The proof uses the local theory, physical estimates and same-witness QFB
tail stated below. They are proved or imported endpoints of the accepted
formal development, rather than unspecified analytic premises of its root
declaration. The budget \eqref{eq:mainbudget}, by contrast, is an explicit
hypothesis of both the mathematical and formal statements.

\begin{remark}[Equivalent physical hypothesis]\label{rem:physical}
On the classical interval, Parseval gives
\begin{equation}\label{eq:budgetconversion}
 \int_{(0,L)}Y(u(t))^2\,dt=16\pi^4\Budget_L(u).
\end{equation}
Consequently \eqref{eq:mainbudget} is equivalently the finiteness of
$\int_{(0,L)}\norm{\nabla u(t)}{2}^4\,dt$. Finiteness on each smaller compact
interval alone is not enough: the bound must be uniform as the endpoint is
approached.
\end{remark}

\section{Local existence, uniqueness and the maximal pair}\label{sec:local}
\begin{proposition}[Common local restart time]\label{prop:local}
There exists $c_0>0$ such that, for each $M\ge0$, every $b\in\D$ with
$\norm{b}{H^1}\le M$ has a unique smooth local normalized solution on an
interval of length at least
\begin{equation}\label{eq:restart}
 h_\nu(M)=\frac{c_0\nu^3}{1+M^4}>0.
\end{equation}
The statement is invariant under time translation.
\end{proposition}
\begin{proof}
The periodic $H^1$ local well-posedness theorem of Tao
\cite[Theorem~5.1, pp.~48--49]{tao}, specialized to zero forcing, supplies
existence, uniqueness, persistence of higher Sobolev regularity and
smoothness for smooth data. Its quantitative viscosity-one condition has
the form $T\norm{b}{H^1}^4\le c$.

For a viscosity-$\nu$ solution define
$v(s,x)=\nu^{-1}u(s/\nu,x)$ and $q(s,x)=\nu^{-2}p(s/\nu,x)$.
Then $(v,q)$ has viscosity one and initial datum $b/\nu$. The corresponding
condition is $\tau\norm{b}{H^1}^4\le c_0\nu^3$, after a fixed norm
equivalence is absorbed into $c_0$. Formula~\eqref{eq:restart} satisfies
this condition on the whole ball. This is precisely the local input used
in \cite[Proposition~3.1]{small}.
\end{proof}

The local result is independent of critical smallness and of $\tq$.
In the formal development its analogue is constructed from physical
Galerkin estimates and an actual limit. The numerical constant of that
construction is not identified with the presentation constant $c_0$.
Only a positive lifespan depending on the common $H^1$ bound is used below.

\begin{lemma}[Uniqueness with normalized pressure]\label{lem:unique}
Two smooth initial fronts with the same datum agree in both velocity and
normalized pressure throughout their common interval.
\end{lemma}
\begin{proof}
Let $z=u-v$ be the velocity difference. On each compact common interval,
the usual difference-energy identity and incompressibility give
\[
 \frac12\frac{d}{dt}\norm{z}{2}^2+\nu\norm{\nabla z}{2}^2
 \le \norm{\nabla v}{\infty}\norm{z}{2}^2.
\]
The coefficient is integrable there, and $z(0)=0$, so Gronwall gives $z=0$.
The equations then give $\nabla(p-q)=0$. A periodic scalar function with
zero gradient is spatially constant; its zero mean makes that constant
zero. Exhaustion by compact intervals proves the assertion.
\end{proof}

\begin{proposition}[Actual maximal existence]\label{prop:maximal}
For each $a\in\D$, the local fronts determine a unique maximal pair
$(u,p)$ on $[0,\Tmax)$ with $\Tmax>0$. Every classical initial front with
the same datum has duration at most $\Tmax$ and agrees with this pair on
its interval.
\end{proposition}
\begin{proof}
Take the supremum in $(0,\infty]$ of the durations of all initial fronts.
The set is nonempty by Proposition~\ref{prop:local}. For any time below
the supremum choose a front extending beyond that time. Lemma~\ref{lem:unique}
makes the values independent of this choice. Every compact shorter
interval lies in one such front, hence the assembled pair is smooth and
solves the physical equations there. The supremum property excludes a
proper extension and gives the stated duration bound.
\end{proof}

\section{How critical smallness generates the Fourier budget}\label{sec:smallbudget}
This section reconstructs the mechanism in \cite[Lemmas~4.1--4.2]{small}
on the \emph{already existing} smooth lifespan. Global existence is not
used to obtain the estimate.

Fix admissible positive Sobolev constants $S_6,S_3$ such that
\begin{equation}\label{eq:sobolev}
 \norm{u}{6}\le S_6\sqrt{Y(u)},\qquad
 \norm{Du}{L^3(\mathrm{op})}\le S_3\sqrt{Z(u)}.
\end{equation}
Here $Du$ is the spatial derivative as a linear map, $\mathrm{op}$ is its
operator norm and the Hilbert--Schmidt norm is used in $Y$. Set
\begin{equation}\label{eq:margin}
 P=S_6S_3,\qquad h_0=\norm{a}{\HH{1/2}},\qquad \delta=\nu-Ph_0.
\end{equation}
The constants are independent of the datum. Their values are not being
optimized or identified with the fixed constants of the audited
small-data roots.

\begin{lemma}[Critical interpolation and product bound]\label{lem:critical}
For a smooth mean-zero solenoidal field,
\begin{equation}\label{eq:criticalbounds}
 Y^2\le EZ,\qquad
 \bigl|\pair{\Psol((u\cdot\nabla)u)}{\Lambda u}\bigr|
 \le P\sqrt E\,Z.
\end{equation}
\end{lemma}
\begin{proof}
Writing the summand of $Y$ as a product of the half- and three-half-order
Fourier weights and applying Cauchy--Schwarz gives $Y\le\sqrt{EZ}$.
The Leray projection is self-adjoint and commutes with $\Lambda$.
H\"older with exponents $(6,3,2)$ therefore gives
\[
 \bigl|\pair{\Psol((u\cdot\nabla)u)}{\Lambda u}\bigr|
 \le \norm{u}{6}\norm{Du}{L^3(\mathrm{op})}\norm{\Lambda u}{2}
 \le P Y\sqrt Z\le P\sqrt E\,Z.
\]
\end{proof}

\begin{proposition}[Invariant ball and quartic estimate]\label{prop:quartic}
If $\delta>0$, then for every $0\le b<\Tmax$,
\begin{align}
 E(u(b))+2\delta\int_0^b Z(u(t))\,dt&\le h_0^2,\label{eq:criticalbudget}\\
 E(u(b))^2+4\delta\int_0^bY(u(t))^2\,dt&\le h_0^4.\label{eq:quartic}
\end{align}
\end{proposition}
\begin{proof}
Testing the projected equation with $\Lambda u$ gives
\begin{equation}\label{eq:criticalevolution}
 E'+2(\nu-P\sqrt E)Z\le0.
\end{equation}
Choose $r$ with $h_0<r<\nu/P$. Before any first exit from $E\le r^2$,
the coefficient of $Z$ is positive and $E$ is nonincreasing. Such an exit
is impossible. Hence $E(t)\le h_0^2$ throughout the lifespan and
$E'+2\delta Z\le0$. Integration proves \eqref{eq:criticalbudget}.
Multiplication by $2E\ge0$ and $Y^2\le EZ$ give
\[
 (E^2)'+4\delta EZ\le0,
 \qquad (E^2)'+4\delta Y^2\le0.
\]
Integration yields \eqref{eq:quartic}.
\end{proof}

\begin{corollary}[Explicit Fourier-budget certificate]\label{cor:smallbudget}
If $\delta>0$ and $L$ is as in \eqref{eq:cutoff}, then
\begin{equation}\label{eq:smallcertificate}
 \boxed{\displaystyle
 \Budget_L(u)\le
 \frac{h_0^4}{64\pi^4\delta}
 =\frac{\norm{a}{\HH{1/2}}^4}
 {64\pi^4\bigl(\nu-P\norm{a}{\HH{1/2}}\bigr)}.}
\end{equation}
\end{corollary}
\begin{proof}
For $L=0$ there is nothing to prove. For $L>0$ choose $b_n\uparrow L$,
with $0\le b_n<L\le\Tmax$. Dropping the nonnegative terminal term in
\eqref{eq:quartic} gives a common bound $h_0^4/(4\delta)$ for the
integrals of $Y^2$. Monotone convergence and
\eqref{eq:budgetconversion} prove \eqref{eq:smallcertificate}.
\end{proof}

The denominator requires \emph{strict} smallness. Neither $\delta=0$
nor $\delta<0$ is covered by this certificate. Outside this small-data
branch, the next sections retain budget finiteness as a hypothesis.

\section{The budget gives a uniform physical bound}\label{sec:gronwall}
\begin{lemma}[Physical enstrophy evolution]\label{lem:enstrophy}
There is a fixed $K>0$ such that every smooth solution satisfies
\begin{equation}\label{eq:enstrophy}
 Y'+\nu W\le K\nu^{-3}\norm{u}{6}^4Y.
\end{equation}
Consequently, with $C=KS_6^4$,
\begin{equation}\label{eq:cubic}
 Y'+\nu W\le C\nu^{-3}Y^3.
\end{equation}
No smallness premise is required.
\end{lemma}
\begin{proof}
Testing \eqref{eq:nse} with $-\Delta u$ eliminates pressure and gives
\[
 \tfrac12Y'+\nu W
 \le \norm{u}{6}\norm{Du}{L^3(\mathrm{op})}\norm{\Delta u}{2}.
\]
The periodic interpolation estimate
$\norm{Du}{L^3(\mathrm{op})}\le C_1Y^{1/4}W^{1/4}$ follows by bounding
the operator norm by the Hilbert--Schmidt norm and interpolating the
gradient between $L^2$ and $L^6$. Young's inequality with exponents
$4/3$ and $4$ yields
\[
 C_1\norm{u}{6}Y^{1/4}W^{3/4}
 \le\tfrac\nu2 W+C_2\nu^{-3}\norm{u}{6}^4Y.
\]
After multiplication by two and an enlargement of the fixed constant,
this is \eqref{eq:enstrophy}. The first inequality in
\eqref{eq:sobolev} gives \eqref{eq:cubic}. This is the estimate of
\cite[Lemma~5.1]{small}.
\end{proof}

\begin{proposition}[One bound on the whole open interval]\label{prop:H1}
Let $0<L\le\Tmax$ be finite and suppose
$M_Y=\int_{(0,L)}Y(t)^2\,dt<\infty$. For every $0\le t<L$,
\begin{align}
 Y(t)&\le Y(0)\exp(C\nu^{-3}M_Y),\label{eq:Ybound}\\
 \norm{u(t)}{H^1}^2
 &\le2Y(0)\exp(C\nu^{-3}M_Y).\label{eq:H1bound}
\end{align}
For the Fourier hypothesis, $M_Y=16\pi^4\Budget_L(u)$.
\end{proposition}
\begin{proof}
Fix a compact interval $[0,b]$ with $b<L$. Dropping $\nu W$ in
\eqref{eq:cubic} gives the linear differential inequality
\[
 Y'\le g(t)Y,\qquad g(t)=C\nu^{-3}Y(t)^2.
\]
Its coefficient is continuous on $[0,b]$ and its integral is at most
$C\nu^{-3}M_Y$. Gronwall proves \eqref{eq:Ybound}, with a bound
independent of $b$. There is no division by $Y(0)$, so the zero case is
included. Mean-zero Poincar\'e on the unit torus gives
\[
 \norm{u}{2}^2\le(2\pi)^{-2}Y,
 \qquad \norm{u}{H^1}^2=\norm{u}{2}^2+Y
 \le\bigl(1+(2\pi)^{-2}\bigr)Y\le2Y.
\]
This proves \eqref{eq:H1bound}.
\end{proof}

The exponential depends on the \emph{actual} full budget of the solution.
The result does not replace that budget by the initial $L^2$ energy or by
a truncated Fourier computation. The initial energy enters separately
through the eventual-regularity time $\tq$.

\section{Continuation and strict passage beyond the cutoff}\label{sec:continue}
\begin{proposition}[Continuation without a terminal trace]\label{prop:continue}
Let $(u,p)$ be a smooth initial front on $[0,T)$, where $0<T<\infty$.
If $\int_{(0,T)}Y(t)^2\,dt<\infty$, then the pair has a smooth extension
to $[0,S)$ for some $S>T$, agreeing with both original fields on $[0,T)$.
\end{proposition}
\begin{proof}
Proposition~\ref{prop:H1} gives one finite bound for the squared $H^1$
norm. Let $M$ be its square root, so $\norm{u(t)}{H^1}\le M$ for
$0\le t<T$. Proposition~\ref{prop:local} gives the same positive
restart duration $h_\nu(M)$ from every such slice.

Choose $0<s<T$ so that $T-s<h_\nu(M)$. Since $s$ lies strictly inside the
front, $u(s)$ is a smooth admissible datum. Restart at time $s$ to obtain
a pair on $[s,s+h_\nu(M))$. Lemma~\ref{lem:unique}, translated in time,
identifies both velocity and pressure with the old pair on the overlap.
Use the old pair near zero and the restarted pair after an interior seam.
They agree on an open neighborhood of that seam, so the resulting pair
is smooth and solves the equations on $[0,s+h_\nu(M))$. Its endpoint is
strictly larger than $T$.
\end{proof}

This argument never evaluates $u(T)$ or assumes that such a value has
already been constructed. It uses the uniformly controlled smooth slices
strictly before $T$.

\begin{proposition}[Strict cutoff overlap]\label{prop:strict}
Under the hypotheses of Theorem~\ref{thm:wfb},
\begin{equation}\label{eq:strict}
 \tq(a,\nu)<\Tmax.
\end{equation}
\end{proposition}
\begin{proof}
Suppose $\Tmax\le\tq$. Then $\Tmax$ is finite and positive, and
$L=\Tmax$. Equation~\eqref{eq:budgetconversion} makes the physical
quartic gradient integral finite on the \emph{whole} open maximal
interval. Proposition~\ref{prop:continue} produces a proper extension,
contradicting maximality. Hence \eqref{eq:strict} holds.

If $\tq=0$, the conclusion follows directly from $\Tmax>0$.
If $\Tmax=\infty$, it follows from finiteness of $\tq$.
The contradiction argument includes the equality case $\Tmax=\tq>0$.
\end{proof}

\section{The original weak trajectory and its regular tail}\label{sec:tail}
This section states the additional QFB input used by WFB. It is not a
consequence attributed to the small-$H^{1/2}$ publication. Its exact
formal endpoints are retained in the pinned source package \cite{methlib}
and in the accepted Methboke closure \cite{methboke}.

\begin{proposition}[Inherited QFB tail and comparison]\label{prop:qfb}
For $\nu>0$ and $a\in\D$, a global Leray--Hopf velocity $v$ with initial
datum $a$ exists. Fix one such original weak trajectory, together with
its weak-gradient witness. The following assertions hold.
\begin{enumerate}[label=(\roman*),leftmargin=2em,itemsep=2pt]
\item On every compact classical interval with datum $a$, weak--strong
comparison identifies $v(t)$ with the classical velocity spatially almost
everywhere, at each time in that interval.
\item For every $d>\tq(a,\nu)$ there is a smooth periodic normalized pair
$(w,q)$ on $[d,\infty)$ solving the evolution and incompressibility
equations in \eqref{eq:nse} on that interval, such that
$v(t)=w(t)$ spatially almost everywhere for every $t\ge d$.
\end{enumerate}
\end{proposition}

The existence and tail declarations are, respectively,
\decl{QuasiFeffermanB.Q01_LERAY_HOPF_EXISTENCE} and
\decl{QuasiFeffermanB.Q09_POST_CUTOFF_SMOOTHNESS}.
Finite weak--strong comparison supplies part (i). The all-times
representative statement is relevant: space-time almost-everywhere
agreement alone must not silently be read as pointwise agreement at
every time. The imported predicates and comparison endpoints provide
the required representatives.

\begin{lemma}[Velocity and pressure on an open overlap]\label{lem:overlap}
Suppose $d<\Tmax$, and let $(w,q)$ be the tail from Proposition~\ref{prop:qfb}
of the \emph{same} original weak trajectory compared with $(u,p)$.
Then, for every $d<t<\Tmax$,
\begin{equation}\label{eq:overlap}
 u(t)=w(t),\qquad p(t)=q(t)
\end{equation}
pointwise in space.
\end{lemma}
\begin{proof}
For fixed $t$ choose a finite classical horizon beyond $t$ but below
$\Tmax$. Both velocities agree spatially almost everywhere with $v(t)$.
Their spatial continuity implies pointwise equality. As this holds at
every time in the open overlap, the smooth velocity fields and all
their derivatives agree there. Subtracting the physical equations gives
$\nabla(p-q)=0$. Periodicity and zero pressure means imply $p=q$.
\end{proof}

\begin{proposition}[A single global splice]\label{prop:splice}
If $\tq<\Tmax$, there is a global smooth normalized solution $(U,\Pi)$
with datum $a$ which agrees with $(u,p)$ at every original time
$0\le t<\Tmax$.
\end{proposition}
\begin{proof}
Choose finite times
\begin{equation}\label{eq:splicechoices}
 \tq<d<\eta<T<\Tmax.
\end{equation}
Such times exist by strict overlap, including when $\Tmax=\infty$.
Fix one original weak witness and apply Proposition~\ref{prop:qfb} to
that witness with cutoff parameter $d$. Let $(w,q)$ be the resulting tail.
Define
\[
 (U,\Pi)(t)=
 \begin{cases}
 (u,p)(t),&0\le t<\eta,\\
 (w,q)(t),&t\ge\eta.
 \end{cases}
\]
Lemma~\ref{lem:overlap} gives equality on an open neighborhood of
$\eta$, so all derivatives match across the seam. The equations,
periodicity, divergence condition, means and initial value are preserved.
For every original time beyond the seam and below $\Tmax$, the same
lemma still applies. Thus the splice preserves the entire maximal pair,
not just one chosen compact segment.
\end{proof}

Choosing a new unrelated weak solution for the tail would not by itself
give the overlap required in this proof. The original weak witness is
fixed before comparison and retained through Q09 and the splice.

\section{Proof of WFB and the ten-step derivation}\label{sec:proof}
\begin{proof}[Proof of Theorem~\ref{thm:wfb}]
Proposition~\ref{prop:maximal} gives actual maximal existence and
uniqueness. The finite-budget hypothesis and Proposition~\ref{prop:strict}
give $\tq<\Tmax$. Proposition~\ref{prop:splice} then supplies the
global pair $(U,\Pi)$ and the required agreement. For each finite $S>0$,
the restriction of this pair to $[0,S)$ is an initial classical front.
Maximality therefore gives $S\le\Tmax$. Since $S$ is arbitrary,
$\Tmax=\infty$.
\end{proof}

For clarity, the complete publication-to-WFB route can now be recorded
without merging the budget-generating and budget-using arguments.

\step{1}{Construct the local and maximal classical pair.}
Use the local $H^1$ theorem, smooth persistence and normalized uniqueness
to obtain $(u,p)$ and $\Tmax>0$
(Propositions~\ref{prop:local}--\ref{prop:maximal}).

\step{2}{Identify critical coercivity on the small-data branch.}
For $\delta=\nu-P\norm{a}{\HH{1/2}}>0$, the interpolation and product
estimates imply $E\le h_0^2$ and $E'+2\delta Z\le0$
(Lemma~\ref{lem:critical} and Proposition~\ref{prop:quartic}).

\step{3}{Extract one quartic gradient budget.}
Multiplication by $2E$ gives
$E(b)^2+4\delta\int_0^bY^2\le h_0^4$ for every shorter interval.
The common bound passes to $(0,L)$ by monotone convergence.

\step{4}{Convert the budget into the full Fourier normalization.}
The identity $Y=4\pi^2Q_u$ gives
$\Budget_L\le h_0^4/(64\pi^4\delta)$ for small data.
For general smooth data, retain $\Budget_L<\infty$ as an explicit
hypothesis (Corollary~\ref{cor:smallbudget}).

\step{5}{Use the enstrophy inequality independently of smallness.}
The physical calculation yields $Y'+\nu W\le C\nu^{-3}Y^3$
for every smooth solution (Lemma~\ref{lem:enstrophy}).

\step{6}{Control all pre-endpoint slices by one $H^1$ bound.}
Gronwall with coefficient $C\nu^{-3}Y^2$, followed by mean-zero
Poincar\'e, gives \eqref{eq:H1bound} uniformly on $(0,L)$
(Proposition~\ref{prop:H1}).

\step{7}{Restart without assuming a terminal trace.}
The common local lifespan permits a restart at an interior smooth slice
close enough to a finite endpoint to extend beyond it
(Proposition~\ref{prop:continue}).

\step{8}{Exclude both an early endpoint and cutoff equality.}
If $\Tmax\le\tq$, the budget controls the whole maximal interval and
Step~7 contradicts maximality. Thus $\tq<\Tmax$
(Proposition~\ref{prop:strict}).

\step{9}{Attach the tail of the same original weak trajectory.}
Q01, finite weak--strong comparison and Q09 give an open overlap.
Pointwise velocity and normalized-pressure agreement make the splice
smooth (Section~\ref{sec:tail}). These are additional QFB inputs.

\step{10}{Conclude infinite lifetime and preserve the pair.}
Every finite restriction of the global splice is an initial front.
Maximality forces $\Tmax=\infty$ and gives the full conclusion of
Theorem~\ref{thm:wfb}.

\section{Small-data inclusion and the remaining universal question}\label{sec:classes}
For fixed $\nu>0$ and admissible $P$, define
\begin{align}
 \mathcal S_\nu
 &=\{a\in\D:P\norm{a}{\HH{1/2}}<\nu\},\label{eq:smallclass}\\
 \mathcal W_\nu
 &=\{a\in\D:\Budget_{\min\{\Tmax(a),\tq(a,\nu)\}}(u_a)<\infty\},
 \label{eq:wclass}\\
 \mathcal R_\nu^{\mathrm{sm}}
 &=\{a\in\D:\text{a global smooth normalized NSE pair exists}\}.
 \label{eq:rclass}
\end{align}
The maximal velocity $u_a$ in \eqref{eq:wclass} is well defined by
Proposition~\ref{prop:maximal}, before any assertion of global regularity.

\begin{corollary}[The precise connection to the publication]\label{cor:inclusion}
\begin{equation}\label{eq:inclusions}
 \mathcal S_\nu\subseteq\mathcal W_\nu\subseteq
 \mathcal R_\nu^{\mathrm{sm}}.
\end{equation}
The first inclusion has the explicit certificate
\eqref{eq:smallcertificate}; the second is WFB.
\end{corollary}
\begin{proof}
For $a\in\mathcal S_\nu$, Corollary~\ref{cor:smallbudget} gives a
finite bound with positive denominator, so $a\in\mathcal W_\nu$.
Theorem~\ref{thm:wfb} gives the second inclusion.
\end{proof}

No strictness assertion is made for these inclusions. In particular, the
notation does not by itself exhibit data beyond the small critical ball.
The small-data sector here is the intersection of the publication's
admissible class with smooth data. Its theorem for rough $H^{1/2}$ data
also needs compactness and a strong initial trace; those arguments are
not absorbed into the present WFB statement.

The unresolved universal assertion relevant to this criterion is
\begin{equation}\label{eq:universal}
 \forall a\in\D,\qquad
 \Budget_{\min\{\Tmax(a),\tq(a,\nu)\}}(u_a)<\infty.
\end{equation}
Combining \eqref{eq:universal} with WFB would give global smooth
regularity for all data in $\D$. Neither the paper's small-data estimate
nor the verified conditional root establishes \eqref{eq:universal}.
The role of the finite cutoff is to limit the required control in time;
it does not supply that control automatically.

\appendix
\section{Exact mathematical and formal interfaces}\label{app:semantics}
\subsection{The physical endpoint}
The root namespace is
\decl{MethbokeAddon.WeakFouriermanB}.
The declaration \decl{smooth_datum_maximal_classical_front_exists}
constructs an actual maximal front for every positive viscosity and every
smooth periodic solenoidal mean-zero datum. Its final signature has no
local-existence premise.

The declaration \decl{Weak_Fourierman_B} takes such a maximal front and
the budget predicate
\begin{quote}\small\ttfamily\raggedright
FourierBudgetFinite (fun t => datumFourierCoeff (u t))\\
\quad (min Tmax (ENNReal.ofReal\\
\qquad (cutoff (criticalL2Representative a) nu))).toReal
\end{quote}
and returns $\Tmax=\top$ together with an actual
\decl{GlobalSmoothNSE} pair and pointwise agreement of velocity and
pressure at every nonnegative original time.

The formal $\mathrm{GlobalSmoothNSE}$ structure states smoothness on the
closed forward spacetime domain, actual time and spatial derivatives,
periodicity, the physical equation, zero divergence, the initial value,
zero velocity mean and zero pressure mean. It is not a definition of
regularity through a certificate that already assumes the conclusion.

\subsection{Normalization dictionary}
\begin{center}
\begin{tabular}{p{0.37\linewidth}p{0.54\linewidth}}
\toprule
Mathematical quantity & Formal interpretation\\
\midrule
$\norm{u}{2}^2$ & \decl{physicalEnergy}; no factor $1/2$\\[4pt]
$Y=\norm{\nabla u}{2}^2$ & \decl{physicalDissipation}\\[4pt]
$W=\norm{\Delta u}{2}^2$ & \decl{physicalLaplacianSquare}\\[4pt]
$Q_u=Y/(4\pi^2)$ & Full nonnegative Fourier enstrophy\\[4pt]
$\int Y^2=16\pi^4\Budget_L$ & Physical--Fourier budget transfer\\[4pt]
$z=\sqrt W$ & Scalar dissipation variable in the cubic estimate\\
\bottomrule
\end{tabular}
\end{center}
The scalar variable $z$ is not the paper's $Z=\norm{u}{\HH{3/2}}^2$.
Likewise $E=\norm{u}{\HH{1/2}}^2$ must not be used in place of the
$L^2$ quantity defining $\tq$.

\subsection{Endpoint conventions}
The extended lifetime can be infinite. Converting $\Tmax$ alone to a
real number before taking a minimum would not express the intended
interval. The accepted statement first forms the finite minimum with
the cutoff, and only then takes \decl{ENNReal.toReal}.
The open interval excludes any value at a finite maximal endpoint.
For zero initial data, $\tq=0$ and the budget is zero; positivity of the
local lifespan already gives strict overlap. No division by initial
enstrophy is used in the Gronwall argument.

\clearpage
\section{Publication steps and formal anchors}\label{app:anchors}
The ten steps in Section~\ref{sec:proof} are expositional steps. They are
not identical to the ten implementation obligations in the repository:
one grouping follows the publication's small-data mechanism, the other
organizes the construction of the general conditional endpoint.
The table below uses the \emph{expositional} numbering.

\small
\begin{longtable}{@{}p{0.055\linewidth}p{0.875\linewidth}@{}}
\toprule
Step & Principal source and formal anchor\\
\midrule
\endfirsthead
\toprule Step & Principal source and formal anchor (continued)\\\midrule
\endhead
\bottomrule\endfoot
1 & \cite[Proposition~3.1]{small}; local existence and maximal assembly:
\decl{uniform_H1_local_initial_classical_front_exists},
\decl{initial_classical_front_unique},
\decl{smooth_datum_maximal_classical_front_exists},
\decl{maximal_classical_front_unique}.\\[7pt]
2 & \cite[Lemmas~4.1--4.2]{small}: interpolation, product bound and invariant
critical ball. The arbitrary-admissible-constant presentation is analytic;
the imported small-data development contains fixed-constant counterparts.\\[7pt]
3 & \cite[Lemma~4.2, Eq.~(4.4)]{small}: terminal quartic estimate and the
common bound before globalization. The paper's Appendix~C distinguishes
its Galerkin and physical-limit versions.\\[7pt]
4 & Exact physical--Fourier conversion:
\decl{smoothSigmaZero_dissipation_eq_fourierEnstrophy};
\decl{initial_front_dissipation_square_integrable_of_fourier_budget}.\\[7pt]
5 & \cite[Lemma~5.1]{small}; physical formal evolution and absorption:
\decl{PHYSICAL_NONLINEAR_ABSORPTION};
\decl{INITIAL_FRONT_PHYSICAL_ENSTROPHY_CUBIC};
\decl{initial_front_enstrophy_scalar_continuity}.\\[7pt]
6 & Uniform physical bounds:
\decl{initial_front_uniform_H1_of_dissipation_square};
\decl{maximal_front_uniform_H1_before_budget_cap}.\\[7pt]
7 & Actual restart and extension:
\decl{initial_front_extend_from_restart};
\decl{initial_front_extends_of_dissipation_square};
\decl{maximal_classical_lifetime_eq_top_of_uniform_H1}.\\[7pt]
8 & Strict cutoff passage:
\decl{maximal_classical_lifetime_strictly_exceeds_cutoff};
\decl{maximal_front_exists_strictly_after_cutoff}.\\[7pt]
9 & Original weak witness and one tail:
\decl{maximal_front_original_leray_hopf};
\decl{maximal_front_tail_overlap_of_same_leray_hopf};
\decl{maximal_front_tail_gluing_of_leray_hopf}.
Imports Q01 and Q09 of Proposition~\ref{prop:qfb}.\\[7pt]
10 & Global root:
\decl{maximal_classical_front_global_of_strict_cutoff};
\decl{Weak_Fourierman_B}.\\
\end{longtable}
\normalsize

\begingroup\raggedright
Except for the fully qualified QFB inputs, names in this table are in
the namespace stated in Appendix~\ref{app:semantics}. Their source files
are under \path{Material/MethbokeAddon/WeakFouriermanB/} at the accepted
revision. The root is in \path{WeakFouriermanB.lean}; key modules include
\path{LocalExistence.lean}, \path{PhysicalFourier.lean},
\path{PhysicalEnstrophyEvolution.lean}, \path{PhysicalBudgetContinuation.lean},
\path{PhysicalEndpointContinuation.lean}, \path{StrictCutoff.lean} and
\path{MaximalTailGluing.lean}.\par
\endgroup

The direct root implementation invokes the previously accepted classical
finite-cap theorem \decl{finite_cap_classical_admission}, then uses
maximal-front agreement and the impossibility of a proper extension.
That finite-cap theorem contains the Q01/comparison/Q09/gluing route.
The explicit same-original-witness theorem in
\path{MaximalTailGluing.lean} is a separately audited contextual realization
of Step~9; it is not represented here as a direct import of the root file.

\section{External verification and reproducibility}\label{app:evidence}
\subsection{Accepted source identity}
The exact WFB proof and its dependency closure were accepted at
\begin{quote}\small\raggedright
Repository: \url{https://github.com/randogelist/Methboke}\\
Proof revision: \nolinkurl{fdb515d8bde35ec8021bbc26f360aca8dfa2855c}\\
External run: \url{https://github.com/randogelist/Methboke/actions/runs/44566880192}.
\end{quote}
This manuscript is a later documentation artifact. Its preparation does
not change the accepted proof revision and does not constitute a new
compiler run or theorem-level audit.

\begin{center}
\begin{tabular}{p{0.48\linewidth}p{0.40\linewidth}}
\toprule
Accepted check & Recorded result\\
\midrule
Modules in the verified closure & 365\\
Accepted modules reused & 321\\
Modules freshly compiled & 44\\
Fresh endpoint audit packages & 55\\
Unique named declaration audits & 265, plus the QFB cutout\\
External controls & 183 passed\\
Permitted reported axioms & \texttt{propext}, \texttt{Classical.choice},
\texttt{Quot.sound}\\
\bottomrule
\end{tabular}
\end{center}
The report records successful compilation, exact-type and theorem-level
axiom checks, evidence publication and final acceptance. No admitted
proof placeholder or custom analytic axiom is reported among the audited
declarations. Reuse of compatible compiled modules is distinguished from
fresh execution of the required endpoint audits.

\subsection{Pinned dependencies and retained evidence}
\begin{description}[style=nextline,leftmargin=0pt,labelindent=0pt,itemsep=4pt]
\item[Lean] Version 4.33.0; compiler revision
\decl{d8b18978322de05a8f3dba51ef03cf5461676c17}.
\item[Mathlib] Revision
\decl{db584cd6d46c92f209a44c0f1c829460d327499d}.
\item[Methlib foundation] Revision
\decl{96033dd69c994bec337a87369a6c27256e724980}.
\item[Retained summary] \path{audits/material-fdb515d8/summary.json}.
The directory also retains the 22 new package audit logs and their
provenance record.
\item[External evidence artifact] Artifact 14553501775 of the accepted run;
ZIP SHA-256
\hash{ec7fbafac42924bc814b7a23f7c090b5f476d84d76193b5f4c0a100c3051612d}.
\end{description}

\subsection{What the evidence establishes}
The evidence attaches to the exact declarations, their concrete
definitions, source closure and pinned compiler environment. The WFB
root retains the finite Fourier-budget premise. Its successful audit is
not a proof of universal budget finiteness.

The small-data paper \cite{small} separately identifies its fixed-constant
audited roots and its independent analytic proof. Its Sections~3--7 are
not represented as having been formalized by the WFB run. Likewise the
arbitrary-admissible-constant reconstruction in Section~\ref{sec:smallbudget}
is not a newly audited endpoint. Compiled fixed-constant intermediate
results and freshly audited named roots should not be conflated.

This manuscript was prepared by reading the mathematical sources and
retained external evidence, with nonexecuting inspection of the Lean
source. No local Lean, Lake or kernel audit was run. The document itself
is built with a standard LaTeX engine; the editable source accompanies
the PDF. A future formal replay should use the recorded source pins and
compatible external caches, and repeat the required exact-type and
theorem-level axiom checks.

\begin{thebibliography}{9}
\raggedright
\bibitem{small}
S.~Suchan, \emph{Global Regularity for Small $H^{1/2}$ Data in the Periodic
Three-Dimensional Navier--Stokes Equations: Quantitative Bounds and Lean
Verification}, preprint, version~1, 13 November 2026.
\href{https://doi.org/10.5281/zenodo.24185514}{doi:10.5281/zenodo.24185514}.
The attached PDF has the internal title \emph{Global Navier--Stokes
regularity for small $H^{1/2}$ data on $\T$: Sobolev-product bounds and
Fourier normalization}. Principal anchors: Proposition~3.1,
Lemmas~4.1--4.2 and~5.1, Proposition~5.2, and Appendices~B--C.

\bibitem{tao}
T.~Tao, \emph{Localisation and compactness properties of the Navier--Stokes
global regularity problem}, Analysis \& PDE \textbf{6} (2013), no.~1,
25--107. \href{https://doi.org/10.2140/apde.2013.6.25}{doi:10.2140/apde.2013.6.25}.
Periodic local well-posedness: Theorem~5.1, pp.~48--49.

\bibitem{methlib}
\emph{Methlib: Quasi-Fefferman-B and critical Sobolev foundations}, source
package, revision
\decl{96033dd69c994bec337a87369a6c27256e724980}.
\href{https://github.com/randogelist/Methlib/tree/96033dd69c994bec337a87369a6c27256e724980}{Immutable source tree}.
QFB tail endpoint: \path{Quasi-Fefferman-B/QuasiFeffermanB/TailChain.lean}.
Native pinned copies are included in the accepted Methboke closure.

\bibitem{methboke}
\emph{Methboke: Weak\_Fourierman\_B}, source package, revision
\decl{fdb515d8bde35ec8021bbc26f360aca8dfa2855c}.
\href{https://github.com/randogelist/Methboke/tree/fdb515d8bde35ec8021bbc26f360aca8dfa2855c}{Immutable source tree}.
External cached verification: run 44566880192; retained summary and
fresh named audits as specified in Appendix~\ref{app:evidence}.
\end{thebibliography}
\end{document}
